\documentclass[11pt]{article}
\usepackage[margin=1in]{geometry}
\usepackage{amsmath,amssymb,mathtools}
\usepackage{fontspec,unicode-math}
\setmainfont{Latin Modern Roman}
\setsansfont{Latin Modern Sans}
\setmathfont{Latin Modern Math}
\usepackage{xurl,hyperref}
\hypersetup{colorlinks=true,urlcolor=blue,linkcolor=blue}
\newcommand{\sref}[2]{\hyperlink{source:#1:#2}{[S#2]}}
\newcommand{\cataloguescope}[1]{\paragraph{Recorded mathematical question.}#1}
\begin{document}
\section{A. Julien: Relationship between Complexity and Cohomology — Problem}\label{statement}

\subsection{Definitions and mathematical statement}

\paragraph{Shared notation.}$\mathbb N=\{1,2,\ldots\}$, and $\mathbb Z,\mathbb Q,\mathbb R,\mathbb C$ denote the integers, rationals, reals and complex numbers. Any other conventions are specified locally.


Let $T$ be an aperiodic repetitive tiling of $\mathbb R^d$, $d\geq1$. Aperiodic means it has no nonzero translation period; repetitive means every finite patch recurs with bounded gaps. Let $p(n)$ count radius-$n$ patches up to translation. Its tiling hull $\Omega$ is the closure of all translates of $T$ in the local tiling topology (agreement on increasingly large balls after increasingly small shifts). Cohomology here is rational Čech cohomology, and $\operatorname{rk}H^*(\Omega;\mathbb Q)$ denotes the total dimension of these rational cohomology groups. The estimate $p(n)=O(n^d)$ means $p(n)\leq Cn^d$ for all sufficiently large integers $n$.
\paragraph{Statement recorded in the supplied source.}
Does $p(n)=O(n^d)$ imply $\operatorname{rk}\check H^*(\Omega;\mathbb Q)<\infty$ for every aperiodic repetitive $d$-dimensional tiling? \sref{statement}{2}
\subsection{Short English statement}
Does complexity of at most order volume force finite rational cohomology for a repetitive aperiodic tiling?
\subsection{Sources}
\begin{description}
\item[\hypertarget{source:statement:1}{S1}] ULAM.AI and the UnsolvedMath Contributors, \emph{UnsolvedMath: A Curated Collection of Open Mathematics Problems}, catalogue entry. CC BY 4.0. \url{https://huggingface.co/datasets/ulamai/UnsolvedMath}.
\item[\hypertarget{source:statement:2}{S2}] Faustin Adiceam (compiler), \emph{Open Problems and Conjectures related to the Theory of Mathematical Quasicrystals}, arXiv:1604.06280v2 (2016), Section 2.5, Problem 2.5.1, p. 8. \url{https://arxiv.org/pdf/1604.06280}.
\end{description}
\label{end:statement}
\end{document}
